\section{收敛模式与向量值积分}\label{sec:03-04}

前三节把标量积分从非负简单函数一直推进到 $L^p$ 空间与对偶表示。现在出现两道新问题。第一道来自
极限：坏点集合即使越来越小，也可能承载越来越高的尖峰；逐点收敛、依测度收敛和 $L^1$ 收敛因而
不能相互替代。第二道来自值域：若被积函数取值于无限维 Banach 空间，逐坐标可测并不自动产生一个
可以由简单函数逼近的向量场。本节先找出防止质量集中的条件，再建立向量值积分；最后把标量
Fubini 的截面论证升级到 Bochner 积分。


\subsection{依测度收敛、一致可积与 Vitali 定理}\label{subsec:3-4-1}
\index{依测度收敛}\index{一致可积}\index{Vitali 收敛定理}

在有限测度空间上，逐点误差、坏点集合的大小与平均误差分别对应三种不同的收敛语言。最短的反例
已经把它们之间的裂缝全部暴露出来。

Vitali 定理在概率极限定理中往往比 DCT 灵活：它不要求寻找一支同时支配整列的函数，而是直接
禁止一族函数把质量压进越来越小的集合。这个条件也成为概率论中控制积分极限的基本紧性条件。

\begin{example}[缩窄而升高的尖峰]\label{ex:3-4-1-1}
在 $([0,1],\Borel([0,1]),m)$ 上令
\[
  f_n(x)=n\1_{(0,1/n)}(x).
\]
对每个 $\varepsilon>0$，当 $n>\varepsilon$ 时
\[
 m\{|f_n|>\varepsilon\}=\frac1n\longrightarrow0,
\]
但 $\|f_n\|_1=1$。而且给定任意 $K>0$，只要取整数 $n>K$，便有
\[
 \int_{\{|f_n|>K\}}|f_n|\,\dd m=1.
\]
因此坏集虽在缩小，全部质量却集中在其中。仅仅知道坏集的测度趋于零，还不足以控制积分。
\end{example}

\begin{definition}[依测度收敛]\label{def:3-4-1-1}
设 $f_n,f$ 是有限测度空间 $(X,\mathcal A,\mu)$ 上 a.e. 取有限值的可测函数。若对每个
$\varepsilon>0$，
\[
  \mu\{x:|f_n(x)-f(x)|>\varepsilon\}\longrightarrow0,
\]
则称 $f_n$ \emph{依测度收敛}于 $f$，记作 $f_n\xrightarrow{\mu}f$。
\end{definition}

当 $\mu(X)=1$ 时，同一概念在概率论中称为\emph{依概率收敛}（convergence in probability）。

\begin{definition}[一致可积]\label{def:3-4-1-2}
设 $\mu(X)<\infty$。族 $\mathcal F\subset L^1(\mu)$ 称为\emph{一致可积}，若
\[
  \lim_{K\to\infty}\sup_{f\in\mathcal F}
  \int_{\{|f|>K\}}|f|\,\dd\mu=0.
  \tag{3.4.1}\label{eq:ui-tail-definition}
\]
\end{definition}

这个定义既控制高度，也控制小集合上的质量。下面给出以后反复使用的等价形式。

\begin{lemma}[一致可积的两种口径]\label{lemma:3-4-1-ui-equivalence}
在有限测度空间上，\eqref{eq:ui-tail-definition} 等价于下列两项同时成立：
\begin{enumerate}
  \item $\sup_{f\in\mathcal F}\|f\|_1<\infty$；
  \item 对每个 $\varepsilon>0$，存在 $\delta>0$，使得 $\mu(E)<\delta$ 时
  \[
    \sup_{f\in\mathcal F}\int_E|f|\,\dd\mu<\varepsilon.
  \]
\end{enumerate}
\end{lemma}

\begin{proof}
先假设 \eqref{eq:ui-tail-definition}。取 $K_0$ 使其后的统一尾积分小于 $1$，则
\[
  \|f\|_1
  \le K_0\mu(X)+\int_{\{|f|>K_0\}}|f|\,\dd\mu
  \le K_0\mu(X)+1,
\]
所以第一项成立。给定 $\varepsilon>0$，再取 $K$ 使统一尾积分小于 $\varepsilon/2$，并令
$\delta=\varepsilon/(2K)$。当 $\mu(E)<\delta$ 时，
\[
 \int_E|f|
 \le \int_{E\cap\{|f|\le K\}}|f|+
      \int_{\{|f|>K\}}|f|
 <K\delta+\frac\varepsilon2=\varepsilon.
\]

反过来，设两项成立，并记 $M=\sup_{f\in\mathcal F}\|f\|_1$。给定 $\varepsilon>0$，由第二项
取相应的 $\delta$。对任意 $K>\max\{1,M/\delta\}$，点态不等式
$\1_{\{|f|>K\}}\le |f|/K$ 给出
\[
  \mu\{|f|>K\}\le \frac{\|f\|_1}{K}<\delta
\]
对所有 $f\in\mathcal F$ 同时成立，因而相应尾积分小于 $\varepsilon$。这就是
\eqref{eq:ui-tail-definition}。
\end{proof}

一致 $L^1$ 有界只给出上面第一项，\cref{ex:3-4-1-1} 正说明第二项不能省略。一个较高次矩的
一致控制则足以排除尖峰。

\begin{example}[$L^p$ 有界族一致可积]\label{ex:3-4-1-2}
设 $\mu(X)<\infty$、$p>1$、$\mathcal F\subset L^p(\mu)$，且
$C:=\sup_{f\in\mathcal F}\|f\|_p<\infty$。有限测度嵌入先给出
$\mathcal F\subset L^1(\mu)$。在 $\{|f|>K\}$ 上有
$|f|\le K^{1-p}|f|^p$，故
\[
 \sup_{f\in\mathcal F}\int_{\{|f|>K\}}|f|\,\dd\mu
 \le K^{1-p}C^p\longrightarrow0.
\]
因此 $\mathcal F$ 一致可积。概率空间只是这个计算中最常见的情形。
\end{example}

依测度收敛也不能保证整列在每一点安定下来。

\begin{example}[打字机序列]\label{ex:3-4-1-3}
对 $m\ge0$ 与 $0\le k<2^m$ 令
\[
 I_{m,k}=[k2^{-m},(k+1)2^{-m}),\qquad f_{2^m+k}=\1_{I_{m,k}}
\]
（在端点 $1$ 处统一置零）。若第 $n$ 项位于第 $m$ 层，则
$m\{|f_n|>1/2\}=2^{-m}$；当 $n\to\infty$ 时 $m\to\infty$，所以 $f_n\xrightarrow{m}0$。
可是每个 $x\in[0,1)$ 在每一层恰落入一个 $I_{m,k}$。沿整列观察，函数值在每层出现一次 $1$，
其余各项为 $0$，故整列在任何 $x\in[0,1)$ 都不收敛。依测度收敛的正确逐点结论只能落在子列上。
\end{example}

\begin{definition}[依测度 Cauchy]\label{def:3-4-1-3}
若对每个 $\varepsilon>0$，
\[
 \mu\{|f_n-f_m|>\varepsilon\}\longrightarrow0
 \qquad(m,n\to\infty),
\]
则称 $(f_n)$ 是\emph{依测度 Cauchy 列}。
\end{definition}

先记录逐点收敛与依测度收敛之间的一条有限测度结论。有限性在支配函数 $1$ 可积这一处使用。

\begin{proposition}[几乎处处收敛蕴含依测度收敛]\label{proposition:3-4-1-ae-measure}
若 $\mu(X)<\infty$ 且 $f_n\to f$ a.e.，则 $f_n\xrightarrow{\mu}f$。
\end{proposition}

\begin{proof}
固定 $\varepsilon>0$。指标函数
$\1_{\{|f_n-f|>\varepsilon\}}$ 几乎处处趋于零，并由可积函数 $1$ 支配。标量 DCT 给
\[
 \mu\{|f_n-f|>\varepsilon\}
 =\int_X\1_{\{|f_n-f|>\varepsilon\}}\,\dd\mu\longrightarrow0.
\]
\end{proof}

\begin{proposition}[子列原理]\label{proposition:3-4-1-1}
设 $\mu(X)<\infty$。则 $f_n\xrightarrow{\mu}f$，当且仅当 $(f_n)$ 的每个子列都有一个进一步
子列几乎处处收敛于 $f$。
\end{proposition}

\begin{proof}
先设 $f_n\xrightarrow{\mu}f$，并从任意子列中仍以 $(f_n)$ 记之。递归选取 $n_k$，使
\[
  \mu(E_k)<2^{-k},\qquad
  E_k:=\{|f_{n_k}-f|>2^{-k}\}.
\]
对每个 $N$，次可加性给
\[
 \mu\!\left(\bigcup_{k\ge N}E_k\right)
 \le\sum_{k\ge N}2^{-k}=2^{1-N}.
\]
集合 $\bigcap_N\bigcup_{k\ge N}E_k$ 是递减集合
$\bigcup_{k\ge N}E_k$ 的交；第一个集合测度有限，故测度的上连续性给其测度为零。对交集外的
$x$，存在 $N(x)$ 使 $k\ge N(x)$ 时 $x\notin E_k$，于是
$|f_{n_k}(x)-f(x)|\le2^{-k}\to0$。

反之，若原列不依测度收敛，则存在 $\varepsilon,\delta>0$ 及一个子列 $(f_{n_k})$，使
\[
  \mu\{|f_{n_k}-f|>\varepsilon\}\ge\delta\qquad(k\ge1).
\]
按照假设，该子列有进一步子列 $(f_{n_{k_j}})$ 几乎处处收敛于 $f$。由
\cref{proposition:3-4-1-ae-measure}，它又依测度收敛于 $f$，所以上述测度应趋于零，与其不小于
$\delta$ 矛盾。
\end{proof}

现在可以精确填补尖峰反例留下的缺口。

\begin{theorem}[Vitali 收敛定理]\label{theorem:3-4-1-2}
设 $\mu(X)<\infty$。若 $f_n\xrightarrow{\mu}f$ 且 $\{f_n:n\ge1\}$ 一致可积，则
$f\in L^1(\mu)$ 且
\[
  \|f_n-f\|_1\longrightarrow0.
\]
反之，$L^1$ 收敛蕴含依测度收敛，并且 $\{f_n:n\ge1\}$ 一致可积。
\end{theorem}

\begin{proof}
先证正向。一致可积给出共同的 $L^1$ 上界：由
\cref{lemma:3-4-1-ui-equivalence}，存在 $M<\infty$ 使 $\|f_n\|_1\le M$。子列原理允许从
原列抽取 $f_{n_j}\to f$ a.e. 的子列。Fatou 引理给
\[
  \int_X|f|\,\dd\mu
  \le\liminf_{j\to\infty}\int_X|f_{n_j}|\,\dd\mu\le M,
\]
故 $f\in L^1$。

还需把统一尾控制传给这个新出现的 $f$，这里不能预先调用待证的 $L^1$ 收敛。固定 $K>0$。在
$f_{n_j}(x)\to f(x)$ 的点上，若 $|f(x)|>K$，则最终有 $|f_{n_j}(x)|>K$，因而
\[
 |f(x)|\1_{\{|f(x)|>K\}}
 \le\liminf_{j\to\infty}
 |f_{n_j}(x)|\1_{\{|f_{n_j}(x)|>K\}}.
\]
再次使用 Fatou 引理，得到
\[
 \int_{\{|f|>K\}}|f|\,\dd\mu
 \le\liminf_j\int_{\{|f_{n_j}|>K\}}|f_{n_j}|\,\dd\mu
 \le\sup_n\int_{\{|f_n|>K\}}|f_n|\,\dd\mu.
 \tag{3.4.2}\label{eq:vitali-limit-tail}
\]

给定 $\varepsilon>0$，由一致可积及 \eqref{eq:vitali-limit-tail} 取 $K$，使 $f_n$ 的统一尾积分
以及 $f$ 的尾积分都小于 $\varepsilon/4$。令
\[
 h_n=|f_n-f|\wedge 2K.
\]
$h_n\xrightarrow{\mu}0$。对任意 $0<\eta<2K$，
\[
 \int_Xh_n\,\dd\mu
 \le \eta\mu(X)+2K\mu\{h_n>\eta\}.
 \tag{3.4.3}\label{eq:bounded-in-measure-l1}
\]
先取 $\eta$ 使第一项小于 $\varepsilon/4$，再令 $n$ 足够大使第二项小于 $\varepsilon/4$，便得
$\int h_n<\varepsilon/2$。逐点有
\[
 |f_n-f|
 \le |f_n|\1_{\{|f_n|>K\}}
    +|f|\1_{\{|f|>K\}}+h_n;
\]
这个不等式可逐情形核对：当两端都不超过 $K$ 时 $h_n=|f_n-f|$；当两端都超过 $K$ 时两个尾项
之和不小于 $|f_n-f|$；若只有一端超过 $K$，则 $|f_n-f|\le2K$ 时 $h_n$ 已等于差值，而
$|f_n-f|>2K$ 时，相应尾项加 $2K$ 不小于两端绝对值之和，因而不小于差值。积分后令
$n\to\infty$，上面三项之和小于 $\varepsilon$，得到
$\|f_n-f\|_1\to0$。

现在设 $\|f_n-f\|_1\to0$。由
$\1_{\{|f_n-f|>\varepsilon\}}\le|f_n-f|/\varepsilon$ 直接积分，得到
\[
 \mu\{|f_n-f|>\varepsilon\}
 \le\varepsilon^{-1}\|f_n-f\|_1\longrightarrow0,
\]
所以依测度收敛。为证一致可积，注意在集合 $A_n(K)=\{|f_n|>K\}$ 上，若
$|f|\le K/2$，则 $|f_n-f|>K/2$，从而 $|f_n|\le2|f_n-f|$；若 $|f|>K/2$，则
$|f_n|\le|f_n-f|+|f|$。故
\[
 \int_{A_n(K)}|f_n|
 \le2\|f_n-f\|_1+
      \int_{\{|f|>K/2\}}|f|.
 \tag{3.4.4}\label{eq:l1-convergence-ui-tail}
\]
先取 $N$ 使 $n\ge N$ 时第一项一致很小，再令 $K\to\infty$ 控制第二项。有限多个
$f_1,\dots,f_{N-1}$ 各自的尾积分也随 $K\to\infty$ 趋零，故整族一致可积。
\end{proof}

图 \ref{fig:3-4-1-vitali-split} 展示证明中的三项分工：式
\eqref{eq:bounded-in-measure-l1} 把截断后的主体积分压到零，一致可积则同时控制两端的高度尾部。

\begin{figure}[htbp]
\centering
\begin{tikzpicture}[x=1cm,y=1cm]
  \draw[->,BookInk!70] (-.2,0)--(9.4,0) node[right,book label] {$x$};
  \draw[->,BookInk!70] (0,-.2)--(0,3.6) node[above,book label] {误差大小};
  \draw[BookWarning,very thick] plot[smooth] coordinates
    {(.3,.25) (1,.45) (1.55,3.1) (2.05,.6) (3,.3) (4.2,.55) (5.1,.35)
     (6,.65) (6.6,2.75) (7.1,.45) (8.8,.25)};
  \draw[BookWarm,dashed] (0,1.7)--(9,1.7) node[right,book label] {$2K$};
  \fill[BookWarning!18] (1.27,0) rectangle (1.83,3.08);
  \fill[BookWarning!18] (6.35,0) rectangle (6.87,2.72);
  \node[book warning node] at (4.1,3.1) {高尾部\\由 UI 控制};
  \draw[book arrow] (3.15,2.82)--(1.72,2.42);
  \draw[book arrow] (5.05,2.82)--(6.55,2.25);
  \node[book strong node] at (4.25,1.05) {截断主体 $h_n$\\坏集测度趋零};
  \node[book warm node] at (7.8,.72) {低于 $\eta$ 的部分\\至多 $\eta\mu(X)$};
\end{tikzpicture}
\caption{Vitali 证明的三块误差：高度尾部、截断后的坏集，以及阈值以下的均匀小量。}
\label{fig:3-4-1-vitali-split}
\end{figure}

一致可积还可以由一个超线性代价函数检测；这个判据经常比直接估计尾积分方便。

\begin{proposition}[de la \ValleePoussin{} 判据]\label{proposition:3-4-1-3}
设 $\mu(X)<\infty$ 且 $\mathcal F\subset L^1(\mu)$。若存在凸的单调递增函数
$\Phi:[0,\infty)\to[0,\infty)$，满足 $\Phi(0)=0$、
\[
 \frac{\Phi(t)}t\longrightarrow\infty
 \quad(t\to\infty),
 \qquad
 \sup_{f\in\mathcal F}\int_X\Phi(|f|)\,\dd\mu<\infty,
 \tag{3.4.5}\label{eq:dlvp-assumption}
\]
则 $\mathcal F$ 一致可积。反过来，在有限测度空间上（特别在概率空间上），任一一致可积族都存在满足
\eqref{eq:dlvp-assumption} 的 $\Phi$。
\end{proposition}

\begin{proof}
凸性与 $\Phi(0)=0$ 蕴含 $t\mapsto\Phi(t)/t$ 在 $(0,\infty)$ 上单调不减：若
$0<s<t$，则 $s=(s/t)t+(1-s/t)0$，所以
$\Phi(s)\le(s/t)\Phi(t)$。对充分大的 $K$ 有 $\Phi(K)>0$，并且在 $|f|>K$ 上有
\[
  |f|\le \frac{K}{\Phi(K)}\Phi(|f|).
\]
若 \eqref{eq:dlvp-assumption} 中积分的共同上界为 $C$，则
\[
 \sup_{f\in\mathcal F}\int_{\{|f|>K\}}|f|
 \le C\frac{K}{\Phi(K)}\longrightarrow0.
\]

反向设 $\mu(X)<\infty$ 且 $\mathcal F$ 一致可积。递归选取阈值 $K_j$，使
\[
 K_j\ge j,\qquad K_{j+1}>K_j,
 \qquad
 \sup_{f\in\mathcal F}\int_{\{|f|>K_j\}}|f|\,\dd\mu\le2^{-j}
 \qquad(j\ge1).
\]
定义
\[
 \Phi(t)=\sum_{j=1}^{\infty}(t-K_j)_+,
 \qquad (u)_+=\max\{u,0\}.
 \tag{3.4.6}\label{eq:dlvp-construction}
\]
因为 $K_j\ge j$，对每个固定 $t$，和中只有有限个非零项；每个有限部分和都凸且递增，且对固定
$t$ 最终等于 $\Phi(t)$，所以 $\Phi$ 凸、递增并满足 $\Phi(0)=0$。给定 $L>0$，取整数
$M>2L$。当 $t\ge2K_M$ 时，对 $1\le j\le M$ 有
$(t-K_j)_+\ge t/2$，故 $\Phi(t)/t\ge M/2>L$。这按定义证明
$\Phi(t)/t\to\infty$。最后由
Tonelli 定理及 $(|f|-K_j)_+\le |f|\1_{\{|f|>K_j\}}$，
\[
 \int_X\Phi(|f|)\,\dd\mu
 \le\sum_{j=1}^{\infty}
      \int_{\{|f|>K_j\}}|f|\,\dd\mu
 \le\sum_{j=1}^{\infty}2^{-j}=1,
\]
且上界与 $f$ 无关。
\end{proof}

依测度收敛本身也有完备性。其证明再次把一列坏集测度估计压缩成一个几乎处处命题。

\begin{proposition}[依测度完备性]\label{proposition:3-4-1-4}
在有限测度空间中，每个依测度 Cauchy 列都有依测度极限；极限在 a.e. 等价意义下唯一。
\end{proposition}

\begin{proof}
递归选取 $n_1<n_2<\cdots$，使
\[
 \mu\{|f_{n_{k+1}}-f_{n_k}|>2^{-k}\}<2^{-k}.
 \tag{3.4.7}\label{eq:measure-cauchy-fast-subsequence}
\]
与子列原理证明相同，尾并的测度估计说明：除去零集后，存在 $k_0(x)$，使 $k\ge k_0(x)$ 时
\[
 |f_{n_{k+1}}(x)-f_{n_k}(x)|\le2^{-k}.
\]
因此向量不是问题，此处只是标量级数：
$\sum_k|f_{n_{k+1}}(x)-f_{n_k}(x)|<\infty$，所以 $(f_{n_k}(x))$ 收敛。把极限记为 $f(x)$，
并在例外零集上置零。它是可测函数，且由
\cref{proposition:3-4-1-ae-measure}，$f_{n_k}\xrightarrow{\mu}f$。

给定 $\varepsilon,\delta>0$。先由原列的 Cauchy 性取 $N$，使 $m,n\ge N$ 时
$\mu\{|f_n-f_m|>\varepsilon/2\}<\delta/2$；再取 $k$ 足够大，使 $n_k\ge N$ 且
$\mu\{|f_{n_k}-f|>\varepsilon/2\}<\delta/2$。于是对所有 $n\ge N$，
\[
 \mu\{|f_n-f|>\varepsilon\}
 \le\mu\{|f_n-f_{n_k}|>\varepsilon/2\}
   +\mu\{|f_{n_k}-f|>\varepsilon/2\}<\delta.
\]
故整列依测度收敛于 $f$。若同时依测度收敛于 $g$，则
\[
 \mu\{|f-g|>\varepsilon\}
 \le\mu\{|f-f_n|>\varepsilon/2\}
   +\mu\{|f_n-g|>\varepsilon/2\}\longrightarrow0,
\]
所以 $f=g$ a.e.
\end{proof}

\begin{remark}[无限空间上要把尖峰与逃逸分开]
若 $\mu(X)=\infty$，a.e. 收敛未必蕴含这里定义的全局依测度收敛。例如
$\1_{[n,n+1]}$ 在 $\R$ 上逐点趋零，但每一项的坏集测度均为 $1$。对一个有限测度耗尽
$E_m\uparrow X$，可以改谈每个 $E_m$ 上的局部依测度收敛；然而局部收敛加高度尾控制仍不足以推出
全局 $L^1$ 收敛，因为质量可能逃向无穷远。还需另加空间紧性条件
\[
 \lim_{m\to\infty}\sup_n\int_{X\setminus E_m}|f_n|\,\dd\mu=0.
\]
$n\1_{(0,1/n)}$ 违反高度尾控制，$\1_{[n,n+1]}$ 则违反空间紧性；两种失效机制需要分别控制。
\end{remark}

Vitali 定理可直接用于幂函数的积分收敛。例如在总质量为一的测度空间上，若 $r>0$、可测函数
$X_n\xrightarrow{\mu}X$，且
$\{|X_n|^r:n\ge1\}$ 一致可积，则子列原理表明每个子列都有 a.e. 收敛于 $X$ 的进一步子列；
在该子列上取连续函数 $t\mapsto|t|^r$，再用子列原理的反向，得到
$|X_n|^r\xrightarrow{\mu}|X|^r$，因而
$\int|X_n|^r\to\int|X|^r$。在概率记号中，这给出依概率收敛加一致可积所蕴含的矩收敛。

\begin{exercise}[尖峰与一致 $L^1$ 界]\label{exercise:3-4-1-1}
构造一个非负函数族，使其 $L^1$ 范数一致有界而不一致可积；再用
\cref{lemma:3-4-1-ui-equivalence} 指出失败的是哪一项。
\end{exercise}
\begin{exercise}[$L^p$ 控制]\label{exercise:3-4-1-2}
设 $p>1$。分别用尾积分估计与小测集绝对连续性证明 $L^p$ 有界族一致可积。
\end{exercise}
\begin{exercise}[打字机子列]\label{exercise:3-4-1-3}
从 \cref{ex:3-4-1-3} 中每层恰取一项。证明所得任意这样的子列都几乎处处趋于零，
并特别写出选取每层第一项时的点态极限。
\end{exercise}
\begin{exercise}[局部 Vitali]\label{exercise:3-4-1-4}
设 $E_m\uparrow X$、$\mu(E_m)<\infty$。在局部依测度收敛、统一高度尾控制及上述空间紧性条件下，
证明 $L^1$ 收敛的对应版本。
\end{exercise}

标量函数可在每个点直接比较大小；向量值函数没有正负部，甚至“每个连续线性泛函看起来都可测”
也可能遗漏值域的不可分部分。下一小节从这个失败处开始。


\subsection{强可测性、Pettis 定理与 Bochner 积分}\label{subsec:3-4-2}
\index{强可测}\index{弱可测}\index{Pettis 可测性定理}\index{Bochner 积分}

本小节中的底层测度空间取完成化。这样，在零集上修改向量值函数仍保持可测性，也与积分的 a.e.
等价类约定一致。若从未完备的 $\sigma$-代数出发，以下陈述应理解为：存在一个 a.e. 相等的代表，
在完成化上满足相应性质。

Hilbert 与 Schmidt 研究层势积分方程时，常把 Fourier 系数组织成 $\ell^2$ 中的向量和无穷矩阵；
有限截断之外还需用完备性补上极限。这里先处理一个更基础的问题：参数化核截面何时真能作为
Banach 或 Hilbert 值函数积分。

例如，对 Hilbert 空间中的规范正交序列 $(e_n)$，形式级数
$u(x)=\sum_na_n(x)e_n$ 即使每个坐标 $a_n$ 都可测，也还不能直接写 $\int u$：必须先证明
部分和确实在范数中收敛，所得极限能由有限值简单函数逼近，并且
$\int(\sum_n|a_n|^2)^{1/2}<\infty$。这说明“逐坐标可测”与“向量可积”之间还隔着范数收敛、
强可测性和范数可积三道条件。

\begin{definition}[强可测]\label{def:3-4-2-1}
设 $B$ 为 Banach 空间。函数 $f:X\to B$ 称为\emph{强可测}，若存在取有限多个值的可测简单函数
$s_n:X\to B$，使 $s_n(x)\to f(x)$ a.e.（收敛关于 $B$ 的范数）。
\end{definition}

\begin{definition}[弱可测与几乎可分值]\label{def:3-4-2-2}
若对每个 $\phi\in B^*$，标量函数 $\phi\circ f$ 都可测，则称 $f$ \emph{弱可测}。若存在
可分闭子空间 $B_0\subset B$，使 $f(x)\in B_0$ a.e.，则称 $f$ \emph{几乎可分值}。
\end{definition}

这里需要的 Hahn--Banach 工具只有\cref{corollary:3-3-3-norming} 已经证明的一点范数化形式。
可分性把逐点可选的泛函压缩成一个可数族。

\begin{lemma}[可分子空间的可数范数化族]\label{lemma:3-4-2-1}
若 $B_0\subset B$ 为可分闭子空间，则存在 $\phi_j\in B^*$、$\|\phi_j\|\le1$，使
\[
  \|b\|=\sup_{j\ge1}|\phi_j(b)|\qquad(b\in B_0).
  \tag{3.4.8}\label{eq:countable-norming-family}
\]
\end{lemma}

\begin{proof}
若 $B_0=\{0\}$，取零泛函即可。否则在 $B_0$ 的单位球面取稠密列 $(u_j)$。由一点
Hahn--Banach，对每个 $j$ 存在 $\psi_j\in B_0^*$，满足
$\|\psi_j\|=1$ 与 $\psi_j(u_j)=1$；再以相同范数把 $\psi_j$ 延拓到 $B^*$，记为 $\phi_j$。
于是 $|\phi_j(b)|\le\|b\|$，所以 \eqref{eq:countable-norming-family} 右端不超过左端。

若 $b\ne0$，给定 $\varepsilon>0$，选 $j$ 使
$\|u_j-b/\|b\|\|<\varepsilon$。则
\[
 |\phi_j(b)|
 =\|b\|\,|\psi_j(b/\|b\|)|
 \ge\|b\|(1-\varepsilon).
\]
令 $\varepsilon\downarrow0$ 即得反向不等式。
\end{proof}

\begin{theorem}[Pettis 可测性定理]\label{theorem:3-4-2-1}
Banach 值函数 $f:X\to B$ 强可测，当且仅当它弱可测且几乎可分值。
\end{theorem}

\begin{proof}
先设 $f$ 强可测，并取简单函数 $s_n\to f$ a.e.。所有 $s_n$ 的值合在一起是可数集；令 $B_0$
为其闭线性包，则 $B_0$ 可分，而且在收敛成立的点上 $f(x)\in B_0$。对任意 $\phi\in B^*$，
$\phi(s_n)$ 可测且在共同零集外趋于 $\phi(f)$。测度空间完备，故这个 a.e. 极限的任意零集修改仍
可测，于是 $f$ 弱可测。

反之，设 $f$ 弱可测且几乎取值于可分闭子空间 $B_0$。在例外零集上把 $f$ 改为 $0$；完成性
保证修改后的函数仍弱可测。取 $B_0$ 中的稠密列 $(b_k)_{k\ge1}$。由
\cref{lemma:3-4-2-1}，对每个 $k$，
\[
 x\longmapsto\|f(x)-b_k\|
 =\sup_{j\ge1}|\phi_j(f(x)-b_k)|
\]
是可数个可测函数的上确界，故可测。对 $n\ge1$，令 $s_n(x)$ 为
$b_1,\dots,b_n$ 中离 $f(x)$ 最近的点；若有多个最近点，取指标最小者。精确地，值层
\[
 \{s_n=b_k\}
 =\bigcap_{\ell<k}\{\|f-b_k\|<\|f-b_\ell\|\}
  \cap\bigcap_{k<\ell\le n}\{\|f-b_k\|\le\|f-b_\ell\|\}
\]
可测，所以 $s_n$ 是有限值简单函数。又因 $(b_k)$ 稠密，
\[
 \|s_n(x)-f(x)\|=\min_{1\le k\le n}\|b_k-f(x)\|\longrightarrow0
\]
对每个 $x$ 成立。这就证明了强可测性，且整个构造没有调用可测选择定理。
\end{proof}

弱可测缺少可分值条件时，结论确实失败。

\begin{example}[不可分正交族]\label{ex:3-4-2-2}
取 $X=[0,1]$ 配 Lebesgue 测度，$H=\ell^2([0,1])$，并以 $e_t$ 表示第 $t$ 个标准单位向量。
定义 $f(t)=e_t$。对任意 $\phi\in H^*$，令
$c_t=\phi(e_t)$。对任意有限集 $F\subset[0,1]$，在线性组合的系数中选取适当相位，并用
Cauchy--Schwarz，可得
\[
 \left(\sum_{t\in F}|c_t|^2\right)^{1/2}
 =\left|\phi\!\left(
   \sum_{t\in F}\frac{\overline{c_t}}{(\sum_{s\in F}|c_s|^2)^{1/2}}e_t
   \right)\right|
 \le\|\phi\|,
\]
其中分母为零时不需计算。于是对每个 $n$，集合
$\{t:|c_t|\ge1/n\}$ 必为有限集，否则可选任意大的有限子集使左端超过 $\|\phi\|$。因此
$\{t:c_t\ne0\}$ 可数，而
$t\mapsto\phi(f(t))=c_t$ 除可数集外为零，故可测。这对每个 $\phi\in H^*$ 成立，所以 $f$ 弱可测。

另一方面，每个 $h\in H$ 的支撑都可数，因为
$\{t:|h_t|\ge1/n\}$ 对每个 $n$ 必为有限集，而支撑是这些集合的可数并。任一可分子空间
$H_0\subset H$ 因而都包含在某个可数坐标子空间中：取 $H_0$ 的稠密列，把各向量的可数支撑取并
得到可数集 $J$，再取闭包便有 $H_0\subset\ell^2(J)$。于是除可数个
$t\in J$ 外，$e_t\notin H_0$。所以 $f$ 不几乎可分值，依
\cref{theorem:3-4-2-1} 也不强可测。
\end{example}

\begin{example}[连续映射的强可测性]\label{ex:3-4-2-3}
设 $X$ 是可分度量空间，其测度结构包含 Borel 集，$B$ 是 Banach 空间，且 $f:X\to B$ 连续。
取 $X$ 的可数稠密集 $D$。连续性给
$f(X)\subset\overline{f(D)}$，所以 $f$ 几乎可分值；对每个 $\phi\in B^*$，$\phi\circ f$ 连续，
因而 Borel 可测。Pettis 定理遂说明 $f$ 强可测。这一结论以后可用于连续参数的核截面。
\end{example}

现在进入积分。逐个对偶坐标可积并不保证这些坐标积分来自同一个 $B$ 中向量，所以不能把它当定义。

\begin{definition}[Bochner 可积]\label{def:3-4-2-3}
强可测函数 $f:X\to B$ 若满足
\[
  \int_X\|f(x)\|\,\dd\mu(x)<\infty,
\]
则称 $f$ \emph{Bochner 可积}。若
$s=\sum_{k=1}^m b_k\1_{E_k}$ 的 $E_k$ 两两不交，且每个 $b_k\ne0$ 所在的 $E_k$ 都有有限测度，
忽略零值层后，定义
\[
  \int_Xs\,\dd\mu=\sum_{k:\,b_k\ne0} b_k\mu(E_k).
  \tag{3.4.9}\label{eq:simple-bochner-integral}
\]
一般 $f$ 的积分由下面的 $L^1(B)$ 逼近构造。
\end{definition}

\begin{lemma}[简单函数逼近与 Bochner 积分的构造]\label{lemma:3-4-2-bochner-construction}
每个 Bochner 可积函数 $f$ 都存在可积简单函数 $s_n$，使
\[
  \int_X\|s_n-f\|\,\dd\mu\longrightarrow0.
  \tag{3.4.10}\label{eq:bochner-simple-l1-approximation}
\]
而且 $\int s_n$ 在 $B$ 中收敛，其极限与逼近列无关。
\end{lemma}

\begin{proof}
先处理简单函数。若
$s=\sum_i b_i\1_{E_i}=\sum_jc_j\1_{F_j}$ 是两个不交层表示，先删去两边的零值层。
由可积性，余下每个值层都有有限测度。集合
$E_i\cap F_j$ 构成共同细分；在每个非空细分层上有 $b_i=c_j$。因此
\[
 \sum_i b_i\mu(E_i)
 =\sum_{i,j}b_i\mu(E_i\cap F_j)
 =\sum_{i,j}c_j\mu(E_i\cap F_j)
 =\sum_jc_j\mu(F_j),
\]
故
\eqref{eq:simple-bochner-integral} 与表示无关。对 $s-t$ 的值层作同样的共同细分，并在有限向量和上
使用三角不等式，得到
\[
 \left\|\int_X(s-t)\,\dd\mu\right\|
 \le\int_X\|s-t\|\,\dd\mu
 \tag{3.4.11}\label{eq:simple-bochner-estimate}
\]
对任意可积简单函数 $s,t$ 成立。

令 $f$ Bochner 可积。由 Pettis 定理的正向部分，修改一个零集后可设 $f(X)$ 包含于可分闭子空间
$B_0$。在 $B_0$ 中取稠密列 $(b_k)_{k\ge1}$，并令 $b_1=0$。对 $f(x)\ne0$，令 $q_n(x)$ 为首个满足
\[
  \|q_n(x)-f(x)\|\le2^{-n}\|f(x)\|
  \tag{3.4.12}\label{eq:relative-countable-approximation}
\]
的 $b_k$；当 $f(x)=0$ 时令 $q_n(x)=0$。稠密性保证候选点存在。函数
$x\mapsto\|f(x)-b_k\|$ 可测，例如可由强简单逼近取连续范数的点态极限；“首个”规则因此使
$q_n$ 的每个值层可测。更具体地，若
$C_{n,k}=\{\|f-b_k\|\le2^{-n}\|f\|\}$，则首个指标为 $k$ 的值层是
$C_{n,k}\setminus\bigcup_{j<k}C_{n,j}$（再把 $\{f=0\}$ 归入零值层）。于是 $q_n$ 是可数值可测
函数，并且
\[
 \|q_n-f\|\le2^{-n}\|f\|,
 \qquad \|q_n\|\le(1+2^{-n})\|f\|.
\]

把 $A_{n,k}$ 取为“首个合格指标恰为 $k$”的两两不交层，并把 $\{f=0\}$ 归入 $A_{n,1}$；于是
$q_n=\sum_kb_k\1_{A_{n,k}}$。由非负积分的可数可加性，
\[
 \sum_{k\ge1}\|b_k\|\mu(A_{n,k})=\int_X\|q_n\|\,\dd\mu<\infty.
\]
故可选 $m(n)$，使这一级数在 $k>m(n)$ 的尾和小于 $2^{-n}$。令
\[
 s_n=\sum_{k\le m(n)}b_k\1_{A_{n,k}},
\]
并在其余值层上置零。它是真正的有限值可积简单函数，且
\[
 \int\|s_n-f\|
 \le\int\|s_n-q_n\|+\int\|q_n-f\|
 \le2^{-n}+2^{-n}\int\|f\|\longrightarrow0.
\]

由 \eqref{eq:simple-bochner-estimate}，
\[
 \left\|\int s_n-\int s_m\right\|
 \le\int\|s_n-f\|+\int\|s_m-f\|\longrightarrow0.
\]
$B$ 完备，所以 $\int s_n$ 有极限。若 $(t_n)$ 是另一套满足
\eqref{eq:bochner-simple-l1-approximation} 的逼近，则
\[
 \left\|\int s_n-\int t_m\right\|
 \le\int\|s_n-f\|+\int\|t_m-f\|,
\]
右端在 $n,m\to\infty$ 时趋零，因此两套极限相同。
\end{proof}

以该极限定义 $\int f$。简单函数估计与 Banach 空间的完备性还给出以下基本性质。

\begin{theorem}[Bochner 积分基本估计]\label{theorem:3-4-2-2}
Bochner 积分线性。若 $f$ Bochner 可积，则
\[
 \left\|\int_Xf\,\dd\mu\right\|
 \le\int_X\|f\|\,\dd\mu,
 \qquad
 \phi\!\left(\int_Xf\,\dd\mu\right)
 =\int_X\phi(f(x))\,\dd\mu(x)
 \tag{3.4.13}\label{eq:bochner-basic-estimates}
\]
对每个 $\phi\in B^*$ 成立。
\end{theorem}

\begin{proof}
线性先对简单函数由有限和成立。若 $s_n\to f$、$t_n\to g$ 于 $L^1(B)$，则
$as_n+bt_n\to af+bg$ 于 $L^1(B)$；在简单函数等式中取极限便得一般线性。

取 \eqref{eq:bochner-simple-l1-approximation} 中的 $s_n$。由
\eqref{eq:simple-bochner-estimate}，
\[
 \left\|\int s_n\right\|\le\int\|s_n\|,
 \qquad
 \left|\int\|s_n\|-\int\|f\|\right|
 \le\int\|s_n-f\|\longrightarrow0.
\]
令 $n\to\infty$ 得第一个估计。对简单 $s_n$，有限和直接给
$\phi(\int s_n)=\int\phi(s_n)$，而
\[
 \int|\phi(s_n)-\phi(f)|
 \le\|\phi\|\int\|s_n-f\|\longrightarrow0.
\]
$\phi$ 的连续性允许左端同时取极限，得到第二式。
\end{proof}

\begin{theorem}[Bochner DCT]\label{theorem:3-4-2-3}
设 $f_n:X\to B$ 强可测，$f_n(x)\to f(x)$ 在范数意义下 a.e.，且
$\|f_n(x)\|\le g(x)$ a.e.，其中 $g\in L^1(\mu)$ 非负。则 $f$ Bochner 可积，
\[
 \int_X\|f_n-f\|\,\dd\mu\longrightarrow0,
 \qquad
 \int_Xf_n\,\dd\mu\longrightarrow\int_Xf\,\dd\mu
\]
（后一收敛在 $B$ 的范数中）。
\end{theorem}

\begin{proof}
对每个 $n$，强可测性给一个几乎可分值子空间 $B_n$。把所有 $B_n$ 的可数稠密集取并，再取闭
线性包，得到一个可分闭子空间 $B_0$；在共同零集外，所有 $f_n(x)$ 以及其极限 $f(x)$ 都属于
$B_0$。又对每个 $\phi\in B^*$，$\phi(f_n)\to\phi(f)$ a.e.，所以完成化上 $f$ 弱可测。
Pettis 定理给出 $f$ 的强可测性。由 $\|f\|\le g$，它范数可积；同样，假设
$\|f_n\|\le g$ 与 $f_n$ 的强可测性说明每个 $f_n$ 都 Bochner 可积。

现在 $\|f_n-f\|\to0$ a.e. 且 $\|f_n-f\|\le2g$。标量 DCT 给第一项收敛；基本估计于是给
\[
 \left\|\int f_n-\int f\right\|
 \le\int\|f_n-f\|\longrightarrow0.
\]
\end{proof}

\begin{example}[Hilbert 值级数]\label{ex:3-4-2-1}
设 $(e_n)$ 是 Hilbert 空间 $H$ 中的规范正交序列，标量函数 $a_n$ 可测，并且
\[
 G(x):=\left(\sum_{n=1}^{\infty}|a_n(x)|^2\right)^{1/2}\in L^1(\mu).
\]
对几乎每个 $x$，正交性说明部分和
$u_N(x)=\sum_{n\le N}a_n(x)e_n$ 在 $H$ 中 Cauchy；记其极限为
$u(x)=\sum_na_n(x)e_n$，并在余下的可测零集上令 $u=0$；则 $\|u(x)\|=G(x)$ a.e.。

每个 $u_N$ 都强可测，因为把有限个可测坐标同时截断并量化即可得到有限值简单逼近。极限 $u$ 的
强可测性也可以由一条显式对角列看出：把 $u_N$ 的前 $N$ 个复坐标的实部与虚部分别量化到网格
$N^{-1}\Z$；当 $\|u_N(x)\|>N$ 时把量化值置零。所得 $v_N$ 只有有限多个值且可测。在任意
$G(x)<\infty$ 的点，当 $N>G(x)$ 后截断不再发生，而坐标量化误差至多
$\sqrt{2N}/N$。因此
\[
 \|v_N(x)-u(x)\|
 \le\|v_N(x)-u_N(x)\|+\|u_N(x)-u(x)\|\longrightarrow0.
\]
所以 $u$ 强可测，并因 $\|u\|=G$ 而 Bochner 可积。又
$\|u_N-u\|\to0$ a.e. 且 $\|u_N-u\|\le2G$，Bochner DCT 给
\[
 \int_Xu_N\,\dd\mu\longrightarrow\int_Xu\,\dd\mu.
\]
这为 Hilbert 值 Fourier 级数逐项积分提供了一个完全可核查的入口。
\end{example}

\begin{proposition}[Banach 值 $L^p$ 的完备性]\label{proposition:3-4-2-4}
对 $1\le p<\infty$，强可测函数模 a.e. 相等、赋范
\[
 \|f\|_{L^p(B)}=\left(\int_X\|f(x)\|^p\,\dd\mu(x)\right)^{1/p},
\]
构成 Banach 空间。取本质上确界范数时，$L^\infty(X;B)$ 也完备。
\end{proposition}

\begin{proof}
强可测函数对线性组合封闭；标量 Minkowski 不等式应用于点态范数后给三角不等式，而
$\|f\|_{L^p(B)}=0$ 等价于 $f=0$ a.e.。因此只需证明完备性。

设 $(f_n)$ 在 $L^p(B)$ 中 Cauchy，$p<\infty$。抽取子列 $(f_{n_k})$，使
\[
 \sum_{k=1}^{\infty}\|f_{n_{k+1}}-f_{n_k}\|_{L^p(B)}<\infty.
 \tag{3.4.14}\label{eq:banach-lp-fast-subsequence}
\]
令 $h_k(x)=\|f_{n_{k+1}}(x)-f_{n_k}(x)\|$。标量 Minkowski 不等式给
\[
 \left\|\sum_{k=1}^Nh_k\right\|_p
 \le\sum_{k=1}^N\|h_k\|_p.
\]
右端由 \eqref{eq:banach-lp-fast-subsequence} 一致有界。对左端的 $p$ 次方用 MCT，得到
$H:=\sum_kh_k\in L^p$，特别地 $H(x)<\infty$ a.e.。在这些点，$B$ 的完备性使
$f_{n_k}(x)$ 收敛于某个 $f(x)$。

每个 $f_{n_k}$ 几乎取值于一个可分闭子空间；把这些子空间的可数稠密集取并并闭线性张成，得到
共同可分闭子空间 $B_0$，而 $f$ 也几乎取值其中。作为强可测函数的 a.e. 点态极限，或由 Pettis
定理，$f$ 强可测。对每个 $j$，
\[
 \|f_{n_j}(x)-f(x)\|
 \le\sum_{k\ge j}h_k(x).
\]
对右端的有限部分用 Minkowski，再令部分和增大并用 MCT，得到
\[
 \|f_{n_j}-f\|_{L^p(B)}
 \le\sum_{k\ge j}\|h_k\|_p\longrightarrow0.
\]
原列为 Cauchy。给定 $\varepsilon>0$，先取 $N$，使 $m,n\ge N$ 时
$\|f_n-f_m\|_{L^p(B)}<\varepsilon/2$；再选 $j$ 使 $n_j\ge N$ 且
$\|f_{n_j}-f\|_{L^p(B)}<\varepsilon/2$。于是对每个 $n\ge N$，三角不等式给出
$\|f_n-f\|_{L^p(B)}<\varepsilon$。

若 $p=\infty$，抽取子列使
$\|f_{n_{k+1}}-f_{n_k}\|_\infty\le2^{-k}$。去掉可数个例外零集后，增量在每一点一致可和，
从而在完备空间 $B$ 中定义极限 $f$，并且
$\|f_{n_j}-f\|_\infty\le\sum_{k\ge j}2^{-k}$。强可测性仍由共同可分值空间与 Pettis 定理得到，
而
\[
  \|f(x)\|\le \|f_{n_1}(x)\|+\sum_{k\ge1}2^{-k}
  \quad\text{a.e.},
\]
所以 $f\in L^\infty(X;B)$。给定 $\varepsilon>0$，先由原列的 Cauchy 性取 $N$，使
$m,n\ge N$ 时 $\|f_n-f_m\|_\infty<\varepsilon/2$；再选 $j$ 使 $n_j\ge N$ 且
$\|f_{n_j}-f\|_\infty<\varepsilon/2$。于是对每个 $n\ge N$，三角不等式给出
$\|f_n-f\|_\infty<\varepsilon$，故整列收敛于 $f$。
\end{proof}

\begin{remark}[选读边界：Bochner 积分与向量测度]
Pettis 可测性说明何时能由简单函数逐点逼近，Bochner 积分则积分一个已经给定的强可测函数；二者
都不蕴含任意绝对连续的 $B$-值向量测度具有 Bochner 密度。后一个性质取决于值域是否具有
Radon--Nikodym 性质（RNP）。自反空间具有 RNP，而某些 $L^\infty$ 型空间不具有 RNP。因此，
从向量测度恢复密度是一个额外的值域几何问题，不能由 Pettis 可测性或 Bochner 积分的定义推出。
本章不使用上述两个 RNP 分类结论，也不把它们当作任何后续证明的前提；这里只用它们标出一个必须在
向量测度理论中另行定义和证明的边界。
\end{remark}

\begin{figure}[htbp]
\centering
\begin{tikzpicture}[node distance=14mm and 16mm]
  \node[book strong node] (f) {向量场 $f:X\to B$};
  \node[book node,right=of f] (weak) {所有坐标 $\phi\circ f$\\可测};
  \node[book warm node,below=of weak] (sep) {值域几乎落在\\可分子空间 $B_0$};
  \node[book strong node,right=18mm of weak] (dist) {可数距离函数\\$\|f-b_k\|$};
  \node[book strong node,below=of dist] (simp) {有限最近点逼近\\$s_n\to f$};
  \draw[book accent arrow] (f)--node[above,book label]{标量观察}(weak);
  \draw[book accent arrow] (weak)--(dist);
  \draw[book accent arrow] (sep)--(dist);
  \draw[book accent arrow] (dist)--(simp);
  \draw[book dashed arrow] (simp.west)--node[below,book label]{强可测}(f.south east);
\end{tikzpicture}
\caption{Pettis 定理的可数化机制：弱可测给坐标，可分值把不可数距离检验压成可数族。}
\label{fig:3-4-2-pettis}
\end{figure}

Bochner 积分允许把参数化核截面视作 $L^2$ 或一般 Banach 空间中的向量，也允许在可积支配下把
向量极限移过积分号。随机 \Ito{} 积分虽然最终也是 $L^2$ 极限，却还需要适应性与等距构造，因而
不是 Bochner 积分的直接特例。

\begin{exercise}[简单向量积分]\label{exercise:3-4-2-1}
从共同分割出发，补全简单函数积分的表示无关性，并证明
\eqref{eq:simple-bochner-estimate}。
\end{exercise}
\begin{exercise}[连续映射]\label{exercise:3-4-2-2}
不用直接引用 \cref{ex:3-4-2-3}，从可分度量空间的一组可数基构造连续 Banach 值函数的简单
逼近。
\end{exercise}
\begin{exercise}[不可分值反例]\label{exercise:3-4-2-3}
补全 \cref{ex:3-4-2-2} 中“任一可分子空间包含在可数坐标子空间中”的闭包论证。
\end{exercise}
\begin{exercise}[Banach 值 $L^p$]\label{exercise:3-4-2-4}
在 \cref{proposition:3-4-2-4} 的证明中，逐步说明 MCT、Minkowski 与 Pettis 定理各自承担的
作用，并验证极限的 a.e. 等价类不依赖所选快速子列。
\end{exercise}

最后还要回答：若 $f$ 同时依赖两个变量，把一个变量积分掉后所得 Banach 值函数是否强可测？逐个
对偶坐标应用标量 Fubini 只能识别一个已经存在的向量，不能自动制造这个向量。下一小节必须单独
完成截面可测性与简单逼近两项工作。


\subsection{向量值 Fubini 与 Hilbert 值应用}\label{subsec:3-4-3}
\index{Bochner Fubini}\index{Hilbert 空间}\index{Schur 测试}

设 $(X,\mathcal A,\mu)$ 与 $(Y,\mathcal B,\nu)$ 都是 $\sigma$-有限测度空间，并采用
\cref{subsec:3-2-1} 中的乘积测度约定。标量 Fubini 已经告诉我们如何控制
$\|f(x,y)\|$，但它本身没有说明 $y\mapsto f(x,y)$ 是否强可测，也没有说明
$x\mapsto\int_Yf(x,y)\,\dd\nu(y)$ 是否是强可测向量场。这两个问题都要从简单函数逼近中获得。

\begin{definition}[Hilbert 空间与规范正交族]\label{def:3-4-3-hilbert}
内积空间 $H$ 若关于范数 $\|h\|=\sqrt{\langle h,h\rangle}$ 完备，则称为\emph{Hilbert 空间}。
族 $(e_j)$ 若满足 $\langle e_j,e_k\rangle=\delta_{jk}$，称为\emph{规范正交族}；若其闭线性包为
$H$，则称为\emph{规范正交基}。本书在复 Hilbert 空间中约定内积对第一变量线性。
\end{definition}

$\ell^2$ 与本章已经证明完备的 $L^2(\mu)$ 是基本例子。在 $H=\C^m$ 中，Bochner 积分就是逐坐标
积分，正交恒等式也只是有限平方和；无限维新增的义务是证明坐标极限仍落在 $H$ 中，并证明积分可与
该极限交换。

\begin{example}[Hilbert 值指示函数]\label{ex:3-4-3-1}
若 $E\in\mathcal A$、$\mu(E)<\infty$ 且 $h\in H$，则
$f(x)=\1_E(x)h$ 是可积简单函数，
\[
  \int_Xf\,\dd\mu=\mu(E)h,
  \qquad
  \int_X\|f\|\,\dd\mu=\mu(E)\|h\|.
\]
若 $h\ne0$ 且 $\mu(E)=\infty$，则第二个积分为无穷；此时第一个表达式不是一个“无穷倍向量”，
而是根本没有 Bochner 积分。
\end{example}

\begin{example}[轨道平均]\label{ex:3-4-3-2}
设 $(U_t)_{t\in\R}$ 是一族酉算子，即每个 $U_t:H\to H$ 线性且保持内积，并且对每个固定
$\xi\in H$，映射
$t\mapsto U_t\xi$ 在范数中连续。若 $g\in L^1(\R)$，则
\[
  t\longmapsto g(t)U_t\xi
\]
强可测：连续轨道的像具有可分闭线性包，乘以标量后仍落在该可分闭线性包中；对每个
$\phi\in H^*$，标量函数 $g(t)\phi(U_t\xi)$ 可测，故可用 Pettis 定理。其范数等于
$|g(t)|\|\xi\|$，所以
$\int_\R g(t)U_t\xi\,\dd t$ 是一个确定的 Hilbert 向量。若 $V:H\to H$ 有界线性，则
\cref{proposition:3-4-3-2} 将严格证明
\[
 V\!\left(\int_\R g(t)U_t\xi\,\dd t\right)
 =\int_\R g(t)VU_t\xi\,\dd t.
\]
\end{example}

\begin{example}[总质量为一时的 Bochner 平均]\label{ex:3-4-3-3}
在总质量为一的测度空间 $(\Omega,\mathcal F,\Prob)$ 上，强可测 Banach 值函数 $Z$ 若满足
$\int\|Z\|\,\dd\Prob<\infty$，可记
\[
  \operatorname{Av}_{\Prob}(Z):=\int_\Omega Z(\omega)\,\dd\Prob(\omega)
\]
为其 Bochner 平均；在概率论中也称为 Bochner 期望。不可分值域中，单有 Borel 可测不能替代
强可测。从绝对连续向量测度恢复密度则是 RNP 问题，与给定强可测函数的 Bochner 积分是不同问题。
\end{example}

\begin{definition}[乘积空间上的 Bochner 可积]\label{def:3-4-3-1}
若 $f:X\times Y\to B$ 关于 $\mathcal A\otimes\mathcal B$ 强可测，并且
\[
 \int_{X\times Y}\|f(x,y)\|\,\dd(\mu\otimes\nu)(x,y)<\infty,
 \tag{3.4.15}\label{eq:bochner-product-integrable}
\]
则称 $f$ 在乘积空间上 Bochner 可积。
\end{definition}

\begin{definition}[积分算子]\label{def:3-4-3-2}
给定标量核 $K:X\times Y\to\C$ 与函数 $g:Y\to\C$，形式表达式
\[
 (Tg)(x)=\int_YK(x,y)g(y)\,\dd\nu(y)
 \tag{3.4.16}\label{eq:integral-operator-definition}
\]
称为由 $K$ 生成的\emph{积分算子}。该式是否 a.e. 绝对存在、是否定义有界算子，都必须由核估计
证明，不能只靠记号宣告。
\end{definition}

\begin{remark}[弱等式的识别原则]\label{def:3-4-3-3}
若 $u,v\in B$ 满足 $\phi(u)=\phi(v)$ 对每个 $\phi\in B^*$ 成立，则一点 Hahn--Banach
范数化推论给 $u=v$。在 Hilbert 空间中，可以等价地检验
$\langle u,h\rangle=\langle v,h\rangle$ 对每个 $h\in H$ 成立。这个原则只能识别两个已经存在的
向量；逐坐标积分等式本身不能保证强可测性或 Bochner 积分的存在。
\end{remark}

\begin{theorem}[Bochner Fubini]\label{theorem:3-4-3-1}
设 $\mu,\nu$ $\sigma$-有限，且 $f:X\times Y\to B$ 关于
$\mathcal A\otimes\mathcal B$ 强可测，并满足
\eqref{eq:bochner-product-integrable}。则对几乎每个 $x$，截面
$f_x(y)=f(x,y)$ 强可测且 Bochner 可积；在例外处置零后，
\[
 F(x):=\int_Yf(x,y)\,\dd\nu(y)
\]
强可测并属于 $L^1(X;B)$，而且
\[
 \int_{X\times Y}f\,\dd(\mu\otimes\nu)
 =\int_X\left(\int_Yf(x,y)\,\dd\nu(y)\right)\dd\mu(x).
 \tag{3.4.17}\label{eq:bochner-fubini}
\]
交换 $X,Y$ 后的结论与迭代积分等式也成立。
\end{theorem}

\begin{proof}
先说明完成化问题。若强可测性相对于乘积测度的完成化给出，将第 $n$ 个简单逼近的
两两不交值层记为 $E_{n,1},\ldots,E_{n,r_n}$。对每层选取
$\mathcal A\otimes\mathcal B$-可测集 $A_{n,k}$，使 $E_{n,k}\mathbin\triangle A_{n,k}$ 为乘积零集，
再依次以 $A_{n,k}\setminus\bigcup_{j<k}A_{n,j}$ 取代 $A_{n,k}$。原值层两两不交，所以这个去交操作
只再改变一个零集，且得到真正的可测简单函数 $\widetilde u_n$。将所有值层的对称差与
原逼近的收敛例外集都包含在某个共同的
$\mathcal A\otimes\mathcal B$-可测乘积零集 $N$ 中；在 $N^c$ 上有 $\widetilde u_n=u_n\to f$。
令 $\widetilde f=\lim_n\widetilde u_n$ 于 $N^c$，并在 $N$ 上置零，则 $\widetilde f$ 具有乘积可测的简单逼近，且
$\widetilde f=f$ a.e.。标量 Tonelli 应用于 $\1_N$，
说明几乎每个截面只改变一个 $\nu$-零集。因此只需对具有乘积可测简单逼近的代表证明结论。

取简单函数 $u_n\to f$ a.e. 于 $X\times Y$。乘积零集的截面几乎处处为 $\nu$-零集，而每个
$(u_n)_x$ 都是 $Y$ 上的可测简单函数。因此对几乎每个 $x$，$(u_n)_x\to f_x$ a.e.，从而
$f_x$ 强可测。对非负标量函数 $\|f\|$ 使用 Tonelli，由
\eqref{eq:bochner-product-integrable} 得
\[
 r(x):=\int_Y\|f(x,y)\|\,\dd\nu(y)<\infty
\]
对几乎每个 $x$ 成立，并且 $\int_Xr\,\dd\mu=\int_{X\times Y}\|f\|$。所以这些截面
Bochner 可积。

还需证明 $F$ 强可测。由\cref{lemma:3-4-2-bochner-construction} 先取一列 $L^1(B)$ 简单逼近，
再抽取子列并重新编号，可使
\[
 \sum_{n=1}^{\infty}\int_{X\times Y}\|f-s_n\|\,\dd(\mu\otimes\nu)<\infty.
 \tag{3.4.18}\label{eq:bochner-fubini-fast-approximation}
\]
若这些简单函数的值层只在乘积完成化中可测，就按上述逐层取代并去交的方法换成
$\mathcal A\otimes\mathcal B$-可测代表；每个 $s_n$ 只改变于乘积零集，所以上式保持不变。
删去零值层后，把 $s_n$ 写成
$s_n=\sum_{k=1}^{m_n}b_{n,k}\1_{E_{n,k}}$，其中 $b_{n,k}\ne0$；可积性保证每个
$(\mu\otimes\nu)(E_{n,k})<\infty$。在所有截面测度有限的点定义
\[
 S_n(x)=\int_Ys_n(x,y)\,\dd\nu(y)
       =\sum_{k=1}^{m_n}b_{n,k}\nu((E_{n,k})_x),
 \tag{3.4.19}\label{eq:simple-section-integral}
\]
并在例外点置零。标量截面可测性定理说明
$x\mapsto\nu((E_{n,k})_x)$ 可测。有限个可测标量函数乘固定向量之和是强可测的：把这些标量函数
同时截断并按越来越细的网格量化，就得到有限值简单逼近。因此 $S_n$ 强可测。

令
\[
 \rho_n(x)=\int_Y\|f(x,y)-s_n(x,y)\|\,\dd\nu(y).
\]
标量 Fubini 及 \eqref{eq:bochner-fubini-fast-approximation} 给
\[
 \sum_n\int_X\rho_n\,\dd\mu<\infty.
\]
对非负级数使用 Tonelli，得 $\sum_n\rho_n(x)<\infty$ a.e.，特别地 $\rho_n(x)\to0$。在这些点，
Bochner 基本估计给
\[
 \|S_n(x)-F(x)\|\le\rho_n(x)\longrightarrow0.
\]
所有 $S_n$ 的几乎可分值子空间取可数闭线性包仍可分，且每个 $\phi\in B^*$ 都满足
$\phi(S_n)\to\phi(F)$ a.e.；Pettis 定理因此给出 $F$ 的强可测性。并且
\[
 \int_X\|S_n-F\|\,\dd\mu
 \le\int_X\rho_n\,\dd\mu
 =\int_{X\times Y}\|s_n-f\|\,\dd(\mu\otimes\nu)\longrightarrow0.
 \tag{3.4.20}\label{eq:section-integrals-l1-converge}
\]
同时 $\|F(x)\|\le r(x)$，所以 $F\in L^1(X;B)$。

对简单函数 $s_n$，标量 Cavalieri 公式逐个值层给
\[
 \int_XS_n(x)\,\dd\mu(x)
 =\sum_kb_{n,k}\int_X\nu((E_{n,k})_x)\,\dd\mu(x)
 =\sum_kb_{n,k}(\mu\otimes\nu)(E_{n,k})
 =\int_{X\times Y}s_n.
\]
由 \eqref{eq:section-integrals-l1-converge} 与 Bochner 基本估计，左端趋于 $\int_XF$；由
$s_n\to f$ 于 $L^1(B)$，右端趋于 $\int_{X\times Y}f$。这证明
\eqref{eq:bochner-fubini}。转置函数 $\widetilde f(y,x)=f(x,y)$ 关于
$\mathcal B\otimes\mathcal A$ 强可测，且其范数积分仍等于 \eqref{eq:bochner-product-integrable}
中的积分；把已经证明的等式应用于 $(Y,\nu)\times(X,\mu)$，便得到另一迭代次序。
\end{proof}

图 \ref{fig:3-4-3-bochner-fubini} 中两条路径之所以汇合，不是因为逐坐标写出了同样的形式式子，而是
因为可积简单函数在两条路径上严格相等，且 Bochner 基本估计允许把等式传到 $L^1(B)$ 极限。

\begin{figure}[htbp]
\centering
\begin{tikzpicture}[node distance=16mm and 23mm]
  \node[book strong node] (prod) {$f\in L^1(X\times Y;B)$};
  \node[book strong node,right=of prod] (sect) {$F(x)=\int_Yf(x,y)\,\dd\nu$\\$F\in L^1(X;B)$};
  \node[book warm node,below=of prod] (direct) {$\int_{X\times Y}f$};
  \node[book warm node,below=of sect] (iter) {$\int_XF(x)\,\dd\mu$};
  \draw[book accent arrow] (prod)--node[above,book label]{沿 $Y$ 积分}(sect);
  \draw[book accent arrow] (prod)--node[left,book label]{直接积分}(direct);
  \draw[book accent arrow] (sect)--node[right,book label]{沿 $X$ 积分}(iter);
  \draw[book accent arrow] (direct)--node[below,book label]{简单函数逼近}(iter);
\end{tikzpicture}
\caption{Bochner Fubini：标量 Tonelli 控制范数，简单函数逼近负责向量等式与强可测性。}
\label{fig:3-4-3-bochner-fubini}
\end{figure}

一旦积分向量存在，有界线性算子便可以安全地穿过积分号。

\begin{proposition}[有界算子与积分交换]\label{proposition:3-4-3-2}
若 $T:B\to C$ 有界线性，$f:X\to B$ Bochner 可积，则 $Tf$ Bochner 可积，并且
\[
  T\!\left(\int_Xf\,\dd\mu\right)=\int_XTf\,\dd\mu.
\]
\end{proposition}

\begin{proof}
由 $f$ 的强简单逼近以及 $T$ 的连续性，$Tf$ 强可测；又有
$\|Tf(x)\|\le\|T\|\,\|f(x)\|$ a.e.，故 $Tf$ 范数可积。
若 $s_n$ 是 $f$ 的 $L^1(B)$ 简单逼近，则 $Ts_n$ 是简单函数，且
\[
 \int\|Ts_n-Tf\|\le\|T\|\int\|s_n-f\|\longrightarrow0.
\]
对简单函数，等式由有限和及 $T$ 的线性成立；令 $n\to\infty$，
$T$ 的连续性与 Bochner 基本估计分别允许等式两端取极限。
\end{proof}

\begin{proposition}[Hilbert 内积交换与 $L^2$ 配对]\label{proposition:3-4-3-3}
若 $f:X\to H$ Bochner 可积，则对每个 $h\in H$，
\[
 \left\langle\int_Xf\,\dd\mu,h\right\rangle
 =\int_X\langle f(x),h\rangle\,\dd\mu(x).
 \tag{3.4.21}\label{eq:hilbert-inner-product-integral}
\]
若 $f,g\in L^2(X;H)$，则 $x\mapsto\langle f(x),g(x)\rangle$ 可积，且
\[
 \int_X|\langle f(x),g(x)\rangle|\,\dd\mu(x)
 \le\|f\|_{L^2(H)}\|g\|_{L^2(H)}.
 \tag{3.4.22}\label{eq:hilbert-l2-pairing}
\]
\end{proposition}

\begin{proof}
映射 $T_h(v)=\langle v,h\rangle$ 是范数为 $\|h\|$ 的有界线性泛函（本书内积对第一变量
线性），所以 \eqref{eq:hilbert-inner-product-integral} 是
\cref{proposition:3-4-3-2} 的特例。

取有限值可测简单函数 $s_n\to f$、$t_n\to g$ a.e.，并把两个例外零集合并。
每个 $\langle s_n,t_n\rangle$ 是标量简单函数，因而可测；内积的连续性给出
$\langle s_n(x),t_n(x)\rangle\to\langle f(x),g(x)\rangle$ 于该零集外。完成测度空间上的零集修改仍可测，
故 $\langle f,g\rangle$ 可测。逐点 Cauchy--Schwarz 及标量 \Holder{} 不等式给
\[
 \int|\langle f,g\rangle|
 \le\int\|f\|\,\|g\|
 \le\left(\int\|f\|^2\right)^{1/2}
      \left(\int\|g\|^2\right)^{1/2},
\]
这就是 \eqref{eq:hilbert-l2-pairing}。注意在无限测度空间上，$f\in L^2$ 本身并不保证
$\int f$ 存在；合法的对象是这里的 $L^2$ 配对。
\end{proof}

积分算子的第一个实用估计完全不需要 Fourier 变换。它把每一行、每一列的 $L^1$ 控制组合成
$L^2$ 算子范数。

\begin{theorem}[Schur 型积分算子估计]\label{theorem:3-4-3-4}
设 $K:X\times Y\to\C$ 联合可测，并存在 $0\le A,B<\infty$，使
\[
 \int_Y|K(x,y)|\,\dd\nu(y)\le A\quad\text{对几乎处处的 }x,
 \qquad
 \int_X|K(x,y)|\,\dd\mu(x)\le B\quad\text{对几乎处处的 }y.
 \tag{3.4.23}\label{eq:schur-row-column}
\]
则 \eqref{eq:integral-operator-definition} 先在有有限测度支撑的简单函数上定义，并唯一延拓为有界
线性算子
\[
 T:L^2(Y,\nu)\longrightarrow L^2(X,\mu),
 \qquad \|T\|\le\sqrt{AB}.
\]
对每个 $g\in L^2(Y)$，该延拓仍与 a.e. 绝对收敛的核积分
$\int_YK(x,y)g(y)\,\dd\nu(y)$ 一致。
\end{theorem}

\begin{proof}
有有限测度支撑的简单函数在 $L^2(Y)$ 中稠密；这由 $\sigma$-有限耗尽、截断和标量简单函数逼近
得到。先取这样的 $g$。由非负 Tonelli 及列估计，
\[
 \int_X\int_Y|K(x,y)|\,|g(y)|^2\,\dd\nu(y)\dd\mu(x)
 \le B\|g\|_2^2<\infty.
 \tag{3.4.24}\label{eq:schur-column-tonelli}
\]
故对几乎每个 $x$，
$H_g(x):=\int_Y|K(x,y)||g(y)|^2\,\dd\nu(y)<\infty$。在这些点，加权
Cauchy--Schwarz 给
\[
\int_Y|K(x,y)g(y)|\,\dd\nu(y)
 \le\left(\int_Y|K(x,y)|\,\dd\nu(y)\right)^{1/2}H_g(x)^{1/2}<\infty
\]
以及
\[
 |Tg(x)|^2
 \le A\int_Y|K(x,y)||g(y)|^2\,\dd\nu(y).
\]
函数 $(x,y)\mapsto K(x,y)g(y)$ 联合可测；分别对其实部、虚部的正负部作参数积分，Cavalieri--Tonelli
说明这些扩展实值函数关于 $x$ 可测。因而在上述绝对可积集合上取积分、在其补集置零所定义的
$Tg$ 是可测函数。
积分并使用 \eqref{eq:schur-column-tonelli}，得到
\[
 \|Tg\|_2^2\le AB\|g\|_2^2.
 \tag{3.4.25}\label{eq:schur-dense-bound}
\]
因此 $T$ 唯一连续延拓到全部 $L^2(Y)$，且范数不超过 $\sqrt{AB}$。若 $A=0$ 或 $B=0$，
\eqref{eq:schur-dense-bound} 已给出零算子结论；而相应的行积分或列积分为零也说明 $K=0$
于乘积几乎处处，故一般 $g$ 的核积分也几乎处处为零。以下可假设 $A,B>0$。

剩下的是识别延拓的点态公式。固定 $g\in L^2(Y)$，取上述稠密子空间中的 $g_n$，满足
$\|g_n-g\|_2\le2^{-n}$，并令
\[
 H_n(x)=\int_Y|K(x,y)|\,|g_n(y)-g(y)|^2\,\dd\nu(y).
\]
Tonelli 与列界给
\[
 \sum_{n=1}^{\infty}\int_XH_n\,\dd\mu
 \le B\sum_{n=1}^{\infty}4^{-n}<\infty.
\]
因此 $\sum_nH_n(x)<\infty$ a.e.，特别地 $H_n(x)\to0$。同样由
\eqref{eq:schur-column-tonelli} 对一般 $g$ 的版本，$H_g(x)<\infty$ a.e.。行估计与加权
Cauchy--Schwarz 因而同时给出
\[
 \int_Y|K(x,y)g(y)|\,\dd\nu(y)<\infty,
 \qquad
 \int_Y|K(x,y)||g_n(y)-g(y)|\,\dd\nu(y)
 \le\sqrt{A H_n(x)}\longrightarrow0
 \tag{3.4.26}\label{eq:schur-pointwise-kernel-limit}
\]
对几乎每个 $x$ 成立。于是核积分 $T_0g_n(x)\to T_0g(x)$ a.e.
这个几乎处处类与 $g$ 的代表无关：若 $g$ 只在 $\nu$-零集 $N$ 上被修改，则由
Tonelli 与列估计，
\[
 \int_X\int_N|K(x,y)|\,\dd\nu(y)\dd\mu(x)
 =\int_N\int_X|K(x,y)|\,\dd\mu(x)\dd\nu(y)=0,
\]
所以修改前后的核积分对几乎每个 $x$ 相同。

另一方面，延拓的算子范数界给
\[
 \sum_n\|Tg_n-Tg\|_2
 \le\sqrt{AB}\sum_n2^{-n}<\infty.
\]
令 $U_N=\sum_{n=1}^N|Tg_n-Tg|$。标量 Minkowski 不等式给出
\[
 \|U_N\|_2\le\sum_{n=1}^N\|Tg_n-Tg\|_2
 \le\sqrt{AB}\sum_{n=1}^{\infty}2^{-n}.
\]
由于 $U_N^2\uparrow(\sum_n|Tg_n-Tg|)^2$，MCT 给出这一无穷和的平方积分有限，
故 $\sum_n|Tg_n(x)-Tg(x)|<\infty$ a.e.；因此 $Tg_n(x)\to Tg(x)$ a.e.。在稠密子空间上
$Tg_n=T_0g_n$，结合 \eqref{eq:schur-pointwise-kernel-limit} 与极限唯一性，得到
\[
 Tg(x)=T_0g(x)=\int_YK(x,y)g(y)\,\dd\nu(y)
\]
几乎处处成立。
\end{proof}

\begin{remark}[一个可回算的 Schur 常数]
在 $X=Y=\R$ 上取 $K(x,y)=e^{-|x-y|}$。平移 $z=x-y$ 给
\[
 \int_\R|K(x,y)|\,\dd y=\int_\R e^{-|z|}\,\dd z=2,
 \qquad
 \int_\R|K(x,y)|\,\dd x=2.
\]
所以 $\|T\|_{2\to2}\le2$。这个算子也可写成与 $e^{-|\cdot|}$ 的卷积；上面的界只用核的行、列
积分即可得到，不需要 Fourier 变换。
\end{remark}

\begin{corollary}[绝对可和的向量级数可逐项积分]\label{corollary:3-4-3-5}
若 $f_n:X\to B$ 强可测且
\[
 \sum_{n=1}^{\infty}\int_X\|f_n(x)\|\,\dd\mu(x)<\infty,
\]
则 $\sum_nf_n(x)$ 在 $B$ 中 a.e. 绝对收敛于一个 Bochner 可积函数 $f$，右侧级数
$\sum_n\int_Xf_n$ 在 $B$ 中绝对收敛，并且
\[
 \int_X\sum_{n=1}^{\infty}f_n(x)\,\dd\mu(x)
 =\sum_{n=1}^{\infty}\int_Xf_n(x)\,\dd\mu(x).
 \tag{3.4.27}\label{eq:bochner-series-integration}
\]
\end{corollary}

\begin{proof}
标量 Tonelli 给
\[
 \int_X\sum_n\|f_n(x)\|\,\dd\mu
 =\sum_n\int_X\|f_n(x)\|\,\dd\mu<\infty,
\]
所以 $\sum_n\|f_n(x)\|<\infty$ a.e.。$B$ 的完备性给出
$f(x)=\sum_nf_n(x)$，并在余下的可测零集上令 $f=0$。所有 $f_n$ 的几乎可分值子空间取可数闭线性包，仍为可分；部分和强可测，
故其点态极限由 Pettis 定理强可测。又
$\|f\|\le\sum_n\|f_n\|\in L^1$，所以 $f$ Bochner 可积。

令 $S_N=\sum_{n\le N}f_n$。则
\[
 \int_X\|f-S_N\|\,\dd\mu
 \le\sum_{n>N}\int_X\|f_n\|\,\dd\mu\longrightarrow0.
\]
Bochner 基本估计给 $\int S_N\to\int f$，而积分线性给
$\int S_N=\sum_{n\le N}\int f_n$。此外
$\sum_n\|\int f_n\|\le\sum_n\int\|f_n\|<\infty$，故右侧在 $B$ 中绝对收敛，并得到
\eqref{eq:bochner-series-integration}。
\end{proof}

取 $f_n(x)=a_n(x)e_n$，\cref{corollary:3-4-3-5} 便回收 Hilbert 值级数逐项积分。这里交换的是
由范数绝对可和保证的级数；右侧在 Banach 空间中绝对收敛，因此任意重排仍有相同的和。一般的
正交平方可和级数若要在积分号下交换，还需另行验证相应
$L^1(H)$ 或 $L^2(H)$ 条件。Schur 条件与 $K\in L^2(X\times Y)$ 的 Hilbert--Schmidt 条件也
不同：前者控制核的行、列绝对积分，后者控制核在乘积空间上的平方积分。

\begin{exercise}[Bochner Fubini 的简单层]\label{exercise:3-4-3-1}
设 $b\in B$、$E\in\mathcal A\otimes\mathcal B$，且 $b=0$ 或
$(\mu\otimes\nu)(E)<\infty$。对 $s=b\1_E$ 直接证明 \eqref{eq:bochner-fubini}，
再推广到可积有限值简单函数；明确指出标量
Cavalieri 公式出现的位置。
\end{exercise}
\begin{exercise}[算子交换]\label{exercise:3-4-3-2}
证明若 $T_n:B\to C$ 在算子范数中趋于 $T$，且 $f$ Bochner 可积，则
$\int T_nf\to\int Tf$，并给出定量误差界。
\end{exercise}
\begin{exercise}[Hilbert 配对]\label{exercise:3-4-3-3}
设 $f\in L^2(X;H)$、$h\in H$ 且 $E\in\mathcal A$、$\mu(E)<\infty$。证明
$\int_Ef$ 存在，并估计 $\|\int_Ef\|$；再用 \eqref{eq:hilbert-inner-product-integral} 计算其与
$h$ 的内积。
\end{exercise}
\begin{exercise}[Schur 测试]\label{exercise:3-4-3-4}
对 $K(x,y)=e^{-a|x-y|}$（$a>0$）计算 Schur 上界。再构造一个满足 Schur 条件但不属于
$L^2(\R^2)$ 的平移不变核，以说明 Schur 条件与 Hilbert--Schmidt 条件不同。
\end{exercise}

至此，本章从简单函数的有限和出发，经过单调极限、乘积积分与 Radon--Nikodym 表示，抵达
Banach 值积分和向量 Fubini。下一章将把积分与平移、缩放以及平均化组合起来，研究卷积、逼近恒等族
与最大函数；那里反复出现的核积分如今都有了严格的存在性与换序依据。
